\chapter{Business Intelligence (BI)}
\label{cap:bi-teoria}

\begin{edital}
Conceitos, fundamentos, técnicas e métodos de BI; sistemas de suporte à
decisão (SSD) e gestão de conteúdo; arquitetura e aplicações de data
warehouse com ETL e OLAP; data warehouse e data mining; visualização (BD
individuais e cubos); mapeamento das fontes de dados; arquitetura de BI.
\end{edital}

\section{O que é BI, e por que não basta consultar o banco transacional}

\begin{conceito}[A pergunta que fundou o assunto]
BI é o conjunto de conceitos e ferramentas para transformar \textbf{dado}
em \textbf{informação} e apoiar a \textbf{tomada de decisão} --- a
pirâmide clássica é dado $\to$ informação $\to$ conhecimento $\to$
inteligência/sabedoria, cada nível agregando contexto e interpretação ao
anterior.

A pergunta concreta que justifica o assunto existir: \textit{por que não
rodar o relatório direto no banco do sistema transacional?} Três motivos:

\begin{itemize}
\item \textbf{Concorrência}: uma consulta que varre cinco anos de vendas
segura recursos que o sistema transacional precisa para atender o caixa
\textbf{agora} --- rodar os dois no mesmo banco compete pelo mesmo
recurso.
\item \textbf{História}: o banco transacional guarda o estado
\textbf{atual} --- ele sabe onde o cliente mora \textbf{hoje}, não onde
morava quando fez a compra que se quer analisar.
\item \textbf{Integração} (o motivo mais caro de resolver): os dados
relevantes estão espalhados em vários sistemas que discordam entre si --- a
mesma unidade organizacional com dois códigos diferentes em dois sistemas,
a mesma data em dois formatos.
\end{itemize}

O BI existe para resolver os três de uma vez, e a arquitetura descrita a
seguir é a solução estruturada para isso.
\end{conceito}

\begin{conceito}[Arquitetura de BI: o caminho do dado, ponta a ponta]
\begin{center}
fontes (OLTP, planilhas, externas) $\to$ \textbf{extração} $\to$
\textit{staging} $\to$ \textbf{transformação} $\to$ \textbf{DW} (e seus
\textit{data marts}) $\to$ \textbf{OLAP/cubo} $\to$ visualização
\end{center}

Cada etapa existe por um motivo específico: a \textbf{extração} lê das
fontes sem incomodá-las mais do que o necessário (janela de carga, leitura
incremental do que mudou desde a última vez). A \textbf{\textit{staging
area}} é o pouso intermediário --- dado bruto já fora da fonte, ainda não
tratado --- e existe para que a fonte seja liberada rapidamente e para que
o tratamento possa ser refeito sem precisar consultar a fonte de novo. A
\textbf{transformação} é onde a integração de fato acontece --- é a etapa
cara do processo, não a simples cópia. O \textbf{DW} guarda o resultado:
integrado, histórico e \textbf{não volátil} (não se atualiza um fato
passado; acrescenta-se um novo). \textbf{OLAP e visualização} são
\textbf{consumo} --- nenhum dos dois transforma dado; quem transforma é o
ETL, e essa separação de responsabilidade é a resposta de vários itens de
prova sobre “em que etapa isso acontece”.

O \textbf{\textit{data mart}} é o recorte do DW por assunto ou
departamento (vendas, RH) --- não é uma cópia menor por economia, é a
entrega organizada no vocabulário de quem vai efetivamente usar aquele
recorte.
\end{conceito}

\section{Data Warehouse}

\begin{conceito}[As quatro características de Inmon]
Repositório central de dados \textbf{integrados, históricos e orientados a
assunto}, destinado a análise, não a transação. As quatro características
canônicas (Inmon): \textbf{orientado a assunto} (organizado em torno de
temas de negócio, como “vendas” ou “cliente”, não em torno de processos de
sistema), \textbf{integrado} (dados de fontes diferentes reconciliados num
formato e vocabulário únicos), \textbf{não volátil} (dado carregado não é
alterado nem apagado, só acrescentado), \textbf{variante no tempo} (cada
registro carrega uma dimensão temporal, permitindo reconstruir o estado
passado).
\end{conceito}

\begin{conceito}[Data Warehouse x Data Lake x Lakehouse]
\textbf{DW}: guarda dado \textbf{tratado e modelado}
(\textit{schema-on-write} --- a estrutura é imposta antes de gravar).
\textbf{Data Lake}: guarda dado \textbf{bruto}, estruturado e não
estruturado (\textit{schema-on-read} --- a estrutura só é imposta na hora
de ler). \textbf{Lakehouse}: une a governança do DW com a flexibilidade do
Lake, tipicamente em camadas: \textbf{Bronze} (bruto, como chegou da
fonte) $\to$ \textbf{Silver} (limpo, deduplicado, com tipos corretos)
$\to$ \textbf{Gold} (modelado dimensionalmente, pronto para consumo de
BI). Um projeto de Machine Learning costuma consumir o \textit{Bronze} ou
\textit{Silver} (quer o dado com o máximo de granularidade e liberdade);
um dashboard de BI consome o \textit{Gold}.
\end{conceito}

\begin{cuidado}
Data Lake \textbf{não substitui} o DW, e não guarda só dado estruturado
(guarda os dois tipos). Lakehouse \textbf{combina} os dois --- não é
sinônimo de nenhum isoladamente.
\end{cuidado}

\subsection{Inmon x Kimball: por onde se começa a construir}

\begin{conceito}[Duas escolas, duas ordens de construção, dois riscos opostos]
As duas escolas discordam sobre \textbf{qual peça se constrói primeiro} ---
e essa discordância tem consequência prática direta no cronograma e no
risco do projeto.

\textbf{Inmon (\textit{top-down})}: constrói-se primeiro o \textbf{DW
corporativo}, normalizado (até a 3FN), como fonte única de verdade da
organização inteira; os \textbf{data marts} dimensionais são
\textbf{derivados} dele depois. Integra melhor e resiste melhor a
mudanças, mas o primeiro relatório demora --- a área usuária espera o
modelo corporativo inteiro ficar pronto antes de ver qualquer valor.

\textbf{Kimball (\textit{bottom-up})}: começa-se por um \textbf{data mart
dimensional} de um único processo de negócio (vendas, por exemplo), já em
esquema estrela, e vão se acrescentando outros processos aos poucos.
Entrega valor cedo. A integração entre os marts não vem de um modelo
central --- vem das \textbf{dimensões conformadas}: a mesma dimensão
\texttt{Produto}, com a mesma chave e o mesmo significado, compartilhada
por todos os marts (a \textit{bus architecture}).

O risco de cada abordagem é o espelho do outro: em Inmon, o risco é o
projeto longo que nunca chega ao usuário; em Kimball, o risco é o
\textbf{silo} --- marts que não conversam entre si porque as dimensões não
foram conformadas de propósito. É por isso que “dimensão conformada” não é
um detalhe de vocabulário: é a peça que faz o \textit{bottom-up}
efetivamente funcionar como um todo coerente, e não como ilhas isoladas.
\end{conceito}

\section{ETL: a ponte para o DW}

\begin{conceito}[O que pertence a cada fase]
\textit{Extract $\to$ Transform $\to$ Load}. O objetivo é extrair de
várias fontes, transformar num formato padronizado e consistente, e
carregar no DW. ETL \textbf{não} é “criar dashboards” nem “treinar modelo
de aprendizado de máquina” --- isso vem depois, como consumo do dado já
carregado.

Como quase tudo interessante acontece dentro do \textit{Transform}, o
critério para separar as fases precisa ser mais fino que “transformar é
mudar o dado”:

\begin{center}
\begin{tabular}{@{}p{2.1cm}p{9.6cm}@{}}
\toprule
\textbf{Fase} & \textbf{O que é dela} \\
\midrule
\textbf{Extract} & ler das fontes (banco, arquivo, API), tipicamente só o
que mudou desde a última carga. \textbf{Sem regra de negócio} --- ela lê,
não julga o conteúdo \\
\textbf{Transform} & \textbf{limpeza} (nulos, formatos, valores
inválidos), \textbf{padronização/conformidade} (o “de-para” que faz dois
sistemas falarem a mesma língua), \textbf{deduplicação}, \textbf{junção}
de fontes diferentes, \textbf{derivação} de medidas calculadas, geração da
\textbf{chave \textit{surrogate}} e aplicação da regra de SCD \\
\textbf{Load} & gravar no destino --- \textbf{dimensões antes da fato}
(porque a fato referencia as chaves das dimensões, então elas precisam já
existir), carga completa x incremental, índices e agregados construídos
no fim \\
\bottomrule
\end{tabular}
\end{center}

A regra prática que resolve a maior parte das questões de “a que fase
pertence esta tarefa”: se a tarefa \textbf{decide} alguma coisa sobre o
conteúdo --- padronizar, deduplicar, casar duas fontes, calcular ---, é
\textbf{Transform}. Ler é \textit{Extract}; gravar é \textit{Load}.
\end{conceito}

\section{OLAP e modelagem dimensional}

\begin{conceito}[Cubo: dimensões + medidas]
Um \textbf{cubo} organiza os dados em \textbf{dimensões} (tempo, produto,
local --- os eixos pelos quais se quer olhar o dado) e
\textbf{métricas/fatos} (vendas, quantidade --- os números que se quer
medir).

Três formas de \textbf{armazenar} fisicamente o cubo: \textbf{ROLAP} (o
dado fica em tabelas relacionais comuns, e o cubo é montado sob demanda
por consulta --- escala melhor para grandes volumes), \textbf{MOLAP} (o
cubo é pré-calculado e guardado numa estrutura própria --- consulta mais
rápida, mas carga mais pesada e demorada) e \textbf{HOLAP} (híbrido:
detalhe fica no relacional, agregados mais usados ficam pré-calculados).
\end{conceito}

\begin{cuidado}
O modelo multidimensional serve para \textbf{análise}, não para
transação --- quem serve à transação é o modelo relacional/OLTP. Ao
lembrar de ROLAP x MOLAP: o \textbf{M} de MOLAP é de
\textit{multidimensional}, o cubo já pré-calculado. “Ferramentas de BI
são incompatíveis com tabelas normalizadas” é falso --- desnormalizar é
uma escolha de desempenho, não uma exigência técnica das ferramentas.
\end{cuidado}

\begin{conceito}[As operações OLAP]
\begin{center}
\begin{tabular}{@{}p{3cm}p{8.6cm}@{}}
\toprule
\textbf{Operação} & \textbf{O que faz} \\
\midrule
\textit{Drill-down} & desce ao detalhe (ano $\to$ mês $\to$ dia) \\
\textit{Roll-up} (\textit{drill-up}) & sobe à agregação (dia $\to$ mês
$\to$ ano) \\
\textit{Slice} & fatia \textbf{uma} dimensão (fixa um único valor nela) \\
\textit{Dice} & subcubo --- filtra em \textbf{várias} dimensões e valores
ao mesmo tempo \\
\textit{Pivot} (\textit{rotate}) & gira a visão, trocando quais dimensões
aparecem em linha e em coluna \\
\bottomrule
\end{tabular}
\end{center}
\end{conceito}

\subsection{Star Schema x Snowflake Schema}

\begin{conceito}[A mesma decisão de desnormalizar ou não, aplicada às dimensões]
\textbf{Star Schema}: uma tabela fato central + dimensões
\textbf{desnormalizadas} (cada dimensão numa única tabela larga, mesmo que
isso repita dado). Menos \textit{joins}, consulta mais rápida --- é o
padrão recomendado por Kimball.

\textbf{Snowflake Schema}: as dimensões são \textbf{normalizadas} em
várias tabelas relacionadas entre si (por exemplo, \texttt{Produto}
aponta para \texttt{Categoria}, que aponta para \texttt{Departamento}).
Mais \textit{joins}, mais lento para consulta analítica, mas economiza
espaço e facilita manter a integridade de hierarquias que mudam.
\end{conceito}

\begin{exemplo}[Contando joins: por que o Star ganha em velocidade]
Para responder “vendas do \textbf{departamento} X no trimestre”:

\begin{center}
\begin{tabular}{@{}p{2.5cm}p{9.2cm}@{}}
\toprule
\textbf{Estrela} & \texttt{Produto} já guarda categoria e departamento
como colunas suas (desnormalizada): \textbf{2 joins} --- fato $\to$
Produto, fato $\to$ Tempo \\
\textbf{Floco de neve} & \texttt{Produto} $\to$ \texttt{Categoria} $\to$
\texttt{Departamento}, cada uma em sua própria tabela: \textbf{4 joins}
para chegar à mesma resposta \\
\bottomrule
\end{tabular}
\end{center}

Menos \textit{joins} $\to$ consulta mais rápida $\to$ é por isso que o
esquema estrela é o padrão recomendado para BI. O floco de neve troca essa
velocidade por economia de espaço e integridade de hierarquia --- uma
escolha legítima em certos contextos, só não é o padrão default.
\end{exemplo}

\begin{cuidado}
Não inverta os nomes: “estrela” com dimensão normalizada em várias
tabelas é, na verdade, floco de neve. Snowflake \textbf{não} é o padrão
recomendado por Kimball --- é o Star.
\end{cuidado}

\subsection{Montando um esquema estrela: as quatro decisões, em ordem}

\begin{conceito}[O projeto dimensional tem uma ordem fixa]
\begin{enumerate}
\item \textbf{Escolher o processo de negócio}: por exemplo, a venda no
varejo.
\item \textbf{Declarar o grão}: “\textbf{uma linha por item de um
pedido}” --- é a decisão que governa todo o resto do projeto (ver
Seção~\ref{sec:granularidade}).
\item \textbf{Escolher as dimensões}: tudo que responde a “por qual
recorte eu quero olhar isto?” --- \texttt{Tempo}, \texttt{Produto},
\texttt{Loja}, \texttt{Cliente}.
\item \textbf{Escolher as medidas}: o que é \textbf{numérico e somável}
no grão declarado --- \texttt{quantidade}, \texttt{valor\_unitário},
\texttt{valor\_total}, \texttt{desconto}.
\end{enumerate}

O resultado é uma \textbf{tabela fato} que contém \textbf{só duas
coisas}: as \textbf{chaves estrangeiras} para as dimensões e as
\textbf{medidas numéricas}. Nenhum texto descritivo mora na fato --- o
nome do produto fica na dimensão \texttt{Produto}. A única exceção com
“nome próprio” é a \textbf{dimensão degenerada}: o
\texttt{número\_do\_pedido} fica na fato porque é um identificador que não
tem atributo nenhum ao redor que justifique uma dimensão própria.
\end{conceito}

\subsection{Granularidade: o grão da fato}
\label{sec:granularidade}

\begin{conceito}[O que uma linha da tabela fato representa]
Definir o grão é a \textbf{primeira} decisão do projeto dimensional, e é
\textbf{irreversível} na prática: tudo o mais (dimensões, medidas,
agregações possíveis) é escolhido depois dela, em função dela.

\begin{center}
\begin{tabular}{@{}p{3.4cm}p{7.6cm}@{}}
\toprule
\textbf{Grão} & \textbf{Consequência} \\
\midrule
\textbf{Fino} (detalhado) --- “uma linha por atendimento” & \textbf{amplia}
as análises possíveis (dá para agregar depois em qualquer nível desejado);
volume maior de dados \\
\textbf{Grosso} (agregado) --- “uma linha por unidade por dia” & tabela
menor e consulta mais rápida, mas o detalhe está \textbf{perdido para
sempre} --- não dá para desagregar depois \\
\bottomrule
\end{tabular}
\end{center}

Regra de Kimball: \textbf{declare o grão no nível mais atômico que a fonte
de dados permitir}. Agregados são construídos \textbf{a partir} do
detalhe, como tabelas adicionais complementares --- nunca no lugar do
detalhe.
\end{conceito}

\begin{cuidado}
Escolher o grão mais detalhado \textbf{amplia} as análises possíveis (não
restringe) --- quem restringe é o agregado. Granularidade é definida pelo
que \textbf{uma linha da fato representa}, não pela quantidade de
dimensões. Não confundir \textbf{grão} (decisão do modelo, fixa desde a
construção) com \textbf{nível de agregação de uma consulta} (decisão de
quem consulta, via \textit{drill-up}, variável a cada consulta).
\end{cuidado}

\subsection{Tipos de fato}

\begin{conceito}
\textbf{Fato aditivo}: soma corretamente em \textbf{todas} as dimensões
(ex.: valor de vendas --- somar vendas de vários dias, produtos ou lojas
sempre faz sentido). \textbf{Fato semiaditivo}: \textbf{não} soma em
todas as dimensões, tipicamente não soma no eixo do \textbf{tempo} (ex.:
saldo bancário, nível de estoque --- somar o saldo de hoje com o de ontem
não produz nada útil, mas somar o saldo de várias contas num mesmo dia
faz sentido). \textbf{Fato não aditivo}: razões, percentuais, médias ---
nunca somam de forma útil em nenhuma dimensão. \textbf{Dimensão
degenerada}: atributo de identificação (número do pedido) que fica na
própria tabela fato, sem ganhar uma dimensão própria.
\end{conceito}

\subsection{Dimensões lentamente mutantes (SCD)}

\begin{conceito}[O que fazer quando um atributo de dimensão muda com o tempo]
Um atributo de dimensão muda (o cliente troca de cidade, o produto troca
de categoria). \textbf{SCD} (\textit{Slowly Changing Dimension}) é a
técnica que decide o que fazer com o histórico correspondente:

\begin{center}
\begin{tabular}{@{}p{1.9cm}p{4cm}p{5.2cm}@{}}
\toprule
\textbf{Tipo} & \textbf{O que faz} & \textbf{Efeito no histórico} \\
\midrule
\textbf{Tipo 1} & \textbf{sobrescreve} o valor antigo & histórico
\textbf{perdido}; os fatos passados passam a ser lidos com o valor novo \\
\textbf{Tipo 2} & \textbf{insere um novo registro} versionado (nova chave
\textit{surrogate}, com \texttt{data\_início}/\texttt{data\_fim} ou
\textit{flag} de corrente) & \textbf{preserva o histórico completo} ---
cada fato continua ligado à versão vigente na época em que ocorreu \\
\textbf{Tipo 3} & guarda o \textbf{valor anterior em outra coluna} da
mesma linha & preserva só \textbf{uma} mudança (o “valor anterior”), não a
série toda de mudanças \\
\bottomrule
\end{tabular}
\end{center}

É o \textbf{Tipo 2} que responde ao cenário clássico: “as vendas antigas
devem continuar associadas à cidade em que o cliente morava na época da
compra”. E é ele que justifica a \textbf{chave \textit{surrogate}} da
dimensão (Capítulo \ref{cap:banco-dados-teoria}): como a mesma chave
natural passa a ter várias versões ao longo do tempo, a fato precisa
apontar para uma chave artificial que identifique a \textbf{versão}
específica, não apenas a entidade genérica.
\end{conceito}

\begin{cuidado}
Aplicar Tipo 1 (sobrescrever) a um cenário que pede preservação de
histórico é erro; aplicar Tipo 2 a um cenário de simples \textbf{correção
de erro de digitação} também é erro no sentido oposto --- correção de erro
não é “mudança de negócio a preservar”, é sobrescrita legítima (Tipo 1).
O Tipo 3 preserva histórico só do \textbf{valor imediatamente anterior},
não da série completa --- é o distrator “quase certo” mais comum entre os
três.
\end{cuidado}

\section{Sistemas de suporte à decisão (SSD/DSS)}

\begin{conceito}
Apoiam gestores em problemas \textbf{estruturados, semiestruturados e não
estruturados} --- a característica central é oferecer flexibilidade
suficiente para vários contextos de decisão diferentes, não se limitar a
um único tipo de problema.
\end{conceito}

\section{Mapeamento de fontes de dados}

\begin{conceito}[Boas práticas, referenciadas pelo DAMA-DMBOK]
\textbf{Fazer}: entrevistar \textit{stakeholders}, analisar sistemas
existentes, avaliar a \textbf{qualidade} e adequação do dado antes de
usá-lo, documentar as fontes como parte da governança de dados.
\textbf{Não fazer}: coletar \textbf{tudo sem discriminação}, incluindo
dados inconsistentes ou irrelevantes, só “para maximizar a quantidade” ---
essa é a prática explicitamente \textbf{não} recomendada.
\end{conceito}

\section{Data Mining}

\begin{conceito}[Quatro tarefas, divididas por ter ou não rótulo prévio]
\textbf{Data Mining} é o processo de descobrir padrões, correlações e
conhecimento não trivial em grandes volumes de dados.

\begin{center}
\begin{tabular}{@{}p{3.2cm}p{3.2cm}p{4.4cm}@{}}
\toprule
\textbf{Tarefa} & \textbf{Tem rótulo?} & \textbf{Cenário típico} \\
\midrule
Classificação & \textbf{sim} (supervisionada) & prever se o cliente vai
inadimplir (classes já conhecidas de antemão) \\
Regressão & \textbf{sim} (supervisionada) & prever um \textbf{valor
numérico} (faturamento do próximo mês) \\
Clusterização & \textbf{não} (não supervisionada) & \textbf{descobrir
grupos} de clientes com comportamento parecido, sem rótulo definido
previamente \\
Regras de associação & não & descobrir \textbf{o que é comprado junto}
(análise de cesta de compras) \\
\bottomrule
\end{tabular}
\end{center}

Em regras de associação: \textbf{suporte} é a frequência do conjunto de
itens no total de transações (quantas vezes A e B aparecem juntos, sobre o
total); \textbf{confiança} é a força da implicação A $\to$ B (das vezes
que A aparece, em quantas B também aparece).
\end{conceito}

\begin{cuidado}
Classificação x clusterização é o par mais confundido: se o enunciado diz
“sem rótulos prévios” ou “descobrir grupos desconhecidos”, é
clusterização; se as classes já existem e o modelo aprende a atribuí-las
a novos casos, é classificação. “Produtos comprados juntos” é sempre
\textbf{regra de associação}, nunca clusterização, mesmo que ambas sejam
técnicas não supervisionadas.
\end{cuidado}

\subsection{CRISP-DM: o ciclo de um projeto de mineração}

\begin{conceito}[Seis fases, iterativas, começando pelo negócio]
\begin{center}
\begin{tabular}{@{}p{0.6cm}p{3.6cm}p{6.8cm}@{}}
\toprule
& \textbf{Fase} & \textbf{O que acontece} \\
\midrule
1 & Entendimento do \textbf{negócio} & objetivos e critérios de sucesso do
negócio --- é por aqui que o ciclo começa, não pelos dados \\
2 & Entendimento dos \textbf{dados} & coleta inicial, exploração,
avaliação de qualidade \\
3 & \textbf{Preparação} dos dados & limpeza, seleção, integração,
formatação \\
4 & \textbf{Modelagem} & escolha da técnica, treino, ajuste de parâmetros
\\
5 & \textbf{Avaliação} & os resultados atendem aos \textbf{objetivos de
negócio} definidos na fase 1? \\
6 & \textbf{Implantação} & colocar o modelo em operação e monitorá-lo \\
\bottomrule
\end{tabular}
\end{center}

O ciclo é \textbf{iterativo}, não uma cascata rígida: volta-se de
qualquer fase para uma anterior conforme necessário, e a seta entre
implantação e entendimento do negócio fecha o círculo, reiniciando o
ciclo com o aprendizado acumulado.
\end{conceito}

\begin{cuidado}
Não confundir a \textbf{avaliação} (fase 5, mede contra os objetivos de
\textbf{negócio}) com a validação técnica do modelo (métricas
estatísticas), que é parte da \textbf{modelagem} (fase 4). O ciclo não
começa pelo entendimento \textbf{dos dados} --- começa pelo
\textbf{negócio}.
\end{cuidado}

\begin{conceito}[KPI x métrica]
Toda métrica é um número medido; \textbf{KPI} (\textit{Key Performance
Indicator}) é uma métrica \textbf{especificamente atrelada a um objetivo
estratégico} da organização --- nem toda métrica é KPI, mas todo KPI é uma
métrica.
\end{conceito}
